\section{Lesser \& Greater Mysteries}

Boris Mouravieff's Gnosis\footnote{\url{http://www.bibliotecapleyades.net/archivos\_pdf/mouravieff3.pdf}} is a bridge between the Lesser Mysteries of orthodox Christianity and the Greater Mysteries promised in the Nordic tradition, as is clear from Mouravieff's passage concerning the existence of two types of humanity on earth:

\begin{quotex}
In the first volume of `Gnosis', we had already referred several times to this coexistence of two essentially different races: one of Men, and another of Anthropoids. We must emphasize the fact that from the esoteric point of view the latter term has no derogatory meaning. First constated very long ago, this fact, although it has been distorted because it is generally seen in a false light, was part of the national, social and judicial consciousness of many ancient and modern peoples. One finds its influence in the Indian idea of the Untouchable, the Greek Helot, the Jewish Goy, the medieval European white Bones and black Bones, the German Nazi Untermensch, etc. ... Incidentally, the legend of blue blood does not belong to the domain of pure fantasy. The error is not in the conception of blue blood as a `psychosomatic' phenomenon, but in the naive medieval belief that this so-called aristocratic blood passes automatically from father to son. Readers of `Gnosis' will easily understand the reasons why this attribute can belong only to twice born beings... 
\end{quotex}


Here we see that ``the sons of God'' \& ``daughters of men'' are a type --- hence, male and female are not unalterable indicators of metaphysical status, nor is skin color; still, the physical type indicates or symbolizes a deeper truth, and is not without significance. More importantly, here we see Mouravieff bringing sacred \& secular tradition together --- this is the spirit of the Traditionalist -- to glean and to synthesize, without slavishly or formulistically following the an expected path (something Gurdjieff\footnote{\url{https://www.youtube.com/watch?v=2tRlNdK00HY&feature=related}} also emphasized).

Mouravieff makes this even more clear here:

\begin{quotex}
Let us make an analogy: a navigator knows how to steer his ship toward any point on the seas or oceans without having been there before. To learn how to do this, he studied the science of navigation, which combines several subjects, such as nautical astronomy; in addition, he learned the use of certain instruments, such as the compass, the sextant, the chronometer, the log, the lead, etc., which enable him to find his position. What is more, he has marine charts at his disposal, and a whole library where he can find a detailed description of every corner of the seas, oceans, islands, of the coasts of continents, of the routes of access to all the gulfs, bays, roadsteads, pons, etc. Alt these elements form the method of approaching the solution to each problem he encounters in his work. Strong in his knowledge and savoir-faire, this navigator, when he receives his captain's order, can make his plans and steer his ship to its destination by the shortest way, although he has never been there before.
\end{quotex}


Here we see there are no guarantees, but a method is possible, where there is a will.

\begin{quotex}
The current of `B' influences from the Absolute II was received by Moses, who transmitted them to his people. It was this that made Israel different from the other nations, who lived immersed in `A' influences under the religious domination of tribal gods. Having thus become the chosen people, depositories of the higher revelation, they received the Promise of the Advent of the Christ-Redeemer, of their own Redemption and that of the other nations through their ministry. Now the Jewish conception of Jehovah was originally the tribal god of Judah, but was later recognized by the other tribes, although with reservations. It was thus elevated to the rank of God of Israel, yet in the religious imagination of the Jews it was never raised above the attributes of the Absolute III, even when, much later, the monotheist conception of the one God came to light. This was a question of a relative monotheism, placing at the height of the celestial pyramid a sort of Demiurge, (`craftsman* in Greek), Jehovah, who was precisely established in the consciousness of the Jewish people as the God of Israel. This deviation is important. It even penetrated into Christianity. In the catechisms we find that Jehovah, (the God of Israel or the image of the Absolute III), is confused with God the Father, the Creator of the Universe.
\end{quotex}


This is Gornahoor's interest in parsing the Yahweh of the OT from God the Father, and even Jesus' Father, of the NT. One cannot disentangle the doctrines without separating the persons (at least in theory). This is not to commit to being a Marcionite; Bondarev argues that Yahweh is the face of God at a particular point in time, \& that Michael is his face now.

\begin{quotex}
Solomon's fall consolidated the traditional dualism that, until then, been intermittent, clearly giving it the meaning and form of an esoteric split. David's Christian tradition, neglected, deformed and largely forgotten by the ruling elite of the chosen people, who were preoccupied with political problems, was received and kept chiefly by the simple people --- rarely by intellectuals --- and continued to make its way silently through the centuries. After the split, the psychic and non-spiritual branch of the Tradition developed its own esotericism, second-class (if one may say so), since it was limited within the boundaries of the Absolute Ill's authority. In its turn, this ritualistic esotericism gave birth to a whole science, equally traditional and hermetic, with Solomon as its head --- Solomon, who is sometimes taken for God himself. Connected with the Temple, this Solomonist tradition of initiation was saved after the destruction of the temple by Titus in 70. This was the last calamity, and gave the signal for the dispersal of Israel. Having become occult, it continued to exist, carefully protected from local Christian persecutions, until the Arab conquest of Palestine gave it a refuge. The ruins of the Temple served as a rallying point and as a sacred symbol for its adepts. Thanks to the Crusades, contacts were established when European knights came to the Holy Land and settled there. The legend of the Knights Templar finding Solomon's treasure in the Temple ruins and making it an object of their special initiations, even though allowed by the Pope, surrounded the White-Coats with the mystical halo of a higher occult science as a complement to their Catholic confession, which remained effective. Parallel to this the Jews, dispersed all over the world, brought their mysticism, which was of Solomonist origin, into Western Europe, where it flourished in their ghettos during the Middle Ages.
\end{quotex}


Mouravieff here dismisses Masonism, \& at the same time explains its association with the Jewish movement in the West, as well as secularization. The important point here is that we need to choose David (a prophet) over Solomon (an earthly king). David was \emph{both}.

\begin{quotex}
It is surprising to constate how easily Christian seekers, (or at any rate, those of Christian origin), brush aside the purely Christian Tradition of Moses-Elias-David which Jesus enriched with the New Testament, along with its projection into the future contained in the Gnosis revealed by out Lord after His Resurrection. Too often, seekers of perfectly good faith omit the Gospel and the Epistles, and delve into the Old Testament and the Solomonist psychic tradition. Theoretically, this research is not harmful. \emph{But if we take into account Jesus' principle by which the disciple cannot be greater than the master}, the work of these seekers cannot lead them any higher than a narrow psychic esotericism, limited to the domain of the Absolute III.
\end{quotex}


The important sentence is ``a disciple is not greater than the master'' -- this is why it is critical to avoid problems in esoteric work which are easily avoided, such as choosing lesser masters, lesser mysteries, lesser paths, all of which really equal the lesser self. Disharmonic focus on pet issues or ideologies will obscure the real matter. Look at the Third Reich, if you doubt this. Choose wisely. Where the treasure is, there the heart is also.

\begin{quotex}
Great Eternity therefore appears as the Great Cycle of Manifestation, encompassing the whole scale of subordinate Cycles and relative Eternities, as well as all Times, which are also relative. Thus it contains the Beginning --- the first creative impulse that begins from the Absolute 0---and goes all the way to the End, that is, to the general and absolute Accomplishment. In this it passes down the whole length of the scale of the Macrocosmos, which contains all relative Accomplishments.The Love from the Absolute 0 fills all Manifestation to its uttermost limits, in all directions and all its specificities, under the aegis of the Absolute I and through the person of the Absolute II, after which, enriched by all the experience gained from one end of the scale to the other, including the kingdom of the Absolute HI, it returns to its source in a primitive, unmanifest state at the heart of the Inexpressible. \textbf{Certain teachings consider this End of ends as a General Annihilation. This is an aberration caused by the psychological structure of our intellect, which is incapable of conceiving ideas outside time and space,} although, in scientific speculation and with the aid of mathematical ideas, we can reach a generally accepted conclusion of the relativity of both. It is a question of an abstraction which, pushed to the limit, is out of reach of the imagination to which human beings can claim to have access solely by means of their Personality in its so-called `normal' state, which is an underdeveloped condition.
\end{quotex}


This gives us an answer to worry over the question of ``annihilating one's self in God'', which had so troubled Evola. Mouravieff's primary concern was to re-vivify the strain of David-Elijah-Jesus and make it public, thereby guaranteeing the safety of the ``Great Work'' of the second birth. This effort can be seen as a means of uniting the Greater \& the Lesser Mysteries, as well as explaining where the Masonic initiations and Jewish influences were a sidetrack for the West.

In summary, Mouravieff's work would appear to be an actual blueprint for reconstructing an esoteric tradition (in that he purports to pass down a valid and real Eastern Orthodox version), a tradition which would actually by a complete Cosmology capable of inter-relating both Romanity
\& the Nordic mythos.



\flrightit{Posted on 2011-08-26 by Logres}

\begin{center}* * *\end{center}

\begin{footnotesize}
\begin{sffamily}
\texttt{Cologero on 2011-08-29 at 00:26 said:}

What did you mean by this statement: ``male and female are not inalterable indicators of metaphysical status''?

\hfill

\texttt{Charlotte Cowell on 2011-08-29 at 08:02 said:}

Firstly, this paragraph rings very true thank you: The important sentence is ``a disciple is not greater than the master'' -- this is why it is critical to avoid problems in esoteric work which are easily avoided, such as choosing lesser masters, lesser mysteries, lesser paths, all of which really equal the lesser self. Disharmonic focus on pet issues or ideologies will obscure the real matter. Look at the Third Reich, if you doubt this. Choose wisely. Where the treasure is, there the heart is also.

Secondly, where you talk about the worrying question of `annihilating one's self in God' (apparently this is inevitable at some stage or other), a little phrase springs to mind as the (re)solution: Solve et Coagula. So yes we will be annihilated, but there is also the possibility of reconstruction closer to the image of God. As it says in the Rubbaiyat:

\begin{verse}
Ah Love! could thou and I with Fate conspire\\
To grasp this sorry Scheme of Things entire,\\
Would we not shatter it to bits -- and thenb\\
Re-mould it nearer to the Heart's Desire!
\end{verse}

\hfill

\texttt{logres on 2011-08-29 at 08:38 said:}

Poorly phrased -- I was suggesting that biological sex is significant, but that it does not preclude (in the case of the woman) an experience of the mysteries.


\hfill

\end{sffamily}
\end{footnotesize}

